(1) Conjugation preserves rank, since $geg^{-1}$ has image $g(\mathrm{im}\, e)$. Conversely let $e, e'$ be idempotents of equal rank $r$ on $V = D^m$. Each splits $V$: $V = \mathrm{im}\,e \oplus \ker e$, both summands free (a finitely generated submodule of a free module over a division ring is free, with well-defined dimension); choose $D$-bases of $\mathrm{im}\,e$ and $\ker e$, and likewise for $e'$. The $D$-linear map $g$ carrying the first adapted basis to the second is invertible and satisfies $g e g^{-1} = e'$, both sides being the projection onto $g(\mathrm{im}\,e) = \mathrm{im}\,e'$ along $g(\ker e) = \ker e'$. In the split case, invertible elements are pairs of invertibles and conjugation acts componentwise (inner automorphisms preserve the two ideals, which are the algebra's minimal central summands), so the same argument applies in each slot and no inner automorphism exchanges slots.

(2) is (1) applied to the pair, $\mathrm{rank}(1-e) = m - \mathrm{rank}(e)$.

(3) If $e$ is central its orbit is $\{e\}$. Conversely, if $geg^{-1} = e$ for every invertible $g$, then $e$ commutes with every invertible element, hence with every sum of two of them, hence with all of $A$ by Lemma \claimref{transport.units-span}: $e$ is central. The identification of the singleton-orbit idempotents with the taxonomy's type C, and of the moved ones with type P, is then Definition \claimref{obj.transport.obstruction-types} read against (1).
